\begin{python0}
from solutions import *; clear()
\end{python0}
\sectionthree{Multiple inheritance}

A class can be the child class of more than one parent class.

Suppose we have a \verb!Vehicle! class and a \verb!Ship! class. Then \verb!AmphibiousVehicle! is a child of both \verb!Vehicle! and \verb!Ship! class.

Syntax for multiple inheritance using public inheritance is simply:

\begin{console}
class P0 {};
class P1 {};
class C: public P0, public P1 {}; 
\end{console}

\sidenote{
  \begin{tikzpicture}
  \node[draw,text width=2cm,minimum height=2cm, line width=0.05cm, inner sep=0.3cm, anchor=center] (P0) at (1,1) 
       {\normalsize\texttt{P0}};
  \node[draw,text width=2cm,minimum height=2cm, line width=0.05cm, inner sep=0.3cm, anchor=center] (P1) at (4,1) 
       {\normalsize\texttt{P1}};
  \node[draw,text width=2cm,minimum height=2cm,minimum width=2cm, line width=0.05cm, inner sep=0.3cm, anchor=north] (C) at (2.5,-2) 
       {\normalsize\texttt{C}};
  \node[below=8mm of P0.south] (MILstart){};
  \node[below=8mm of P1.south] (MILend){};
  \node[above=11mm of C.north] (MILmid){};
  \draw[ultra thick] (MILstart) -- (MILend);
  \draw[ultra thick, ->] (MILend) -- (P1.south);
  \draw[ultra thick, <-] (P0.south) -- (MILstart);
  \draw[ultra thick] (MILmid) -- (C.north);
\end{tikzpicture}
}
Some programming languages do not allow multiple inheritance. Why?

Well, suppose class \verb!C0! is derived from \verb!P0!, \verb!P1! and both \verb!P0!, \verb!P1! have \verb!f()!. If \verb!c! is a \verb!C! object, then \verb!c.f()! is ambiguous. MAKE SURE you run the following:
\begin{console}
#include <iostream>

class P0
{
public:
        void f() {}
};

class P1
{
public:
        void f() {}
};

class C: public P0, public P1
{};

int main()
{   
    C c;
    c.f(); // which one!?!?!
    return 0;
}
\end{console}

Your C++ compiler will probably yell at you and say you have an
ambiguous invocation. This is of course the same for ambiguous member
variables:
\begin{console}
#include <iostream>

class P0
{
public:
        int x_;
};

class P1
{
public:
        int x_;
};

class C: public P0, public P1
{};

int main()
{   
    C c;
    c.x_; // which one!?!?!
    return 0;
}
\end{console}

However note that if you have two members in both parents with the same
name but you never call upon that name, your C++ will \textbf{not}
complain. For instance
\begin{console}
#include <iostream>

class P0
{
public:
        int x_;
};

class P1
{
public:
        int x_;
};

class C: public P0, public P1
{};

int main()
{   
    C c;
    return 0;
}
\end{console}

will compile. If \verb!P0! has a method that works with \verb!P0::x_! and \verb!P1! has a method that works with \verb!P1::x_!, that wouldn't be a problem:
\begin{console}
#include <iostream>

class P0
{
public:
        void f0() { std::cout << x_ << '\n'; }
        int x_;
};

class P1
{
public:
        void f1() { std::cout << x_ << '\n'; }
        int x_;
};

class C: public P0, public P1
{};

int main()
{   
    C c;
    c.f0();
    c.f1();
    return 0;
}
\end{console}

Another situation where multiple inheritance might cause a problem is
when, further up, \verb!P0! and \verb!P1! inherits from a common class:
\begin{python}
from latextool_basic import *
p = Plot()
#p += Grid(x0=-3,y0=-5,x1=3,y1=5)

G = Rect(-1, 3, 1, 5, linewidth=0.05, font=r'\huge', label=r'\texttt{G}')
P0 = Rect(-3, -1, -0.5, 1, linewidth=0.05, font=r'\huge', label=r'\texttt{P0}')
P1 = Rect(0.5, -1, 3, 1, linewidth=0.05, font=r'\huge', label=r'\texttt{P1}')
C = Rect(-1, -5, 1, -3, linewidth=0.05, font=r'\huge', label=r'\texttt{C}')

#points_bottom = [(P0.bottom()), (-2, P0.bottomy()), (2, P1.bottomy()), P1.bottom(), (2, P1.bottomy()), (0, -2), (C.top())]

#print(points_bottom)

p += Line(points=[(-2,-1),(-2,-2),(2,-2),(2,-1),(2,-2),(0,-2),(C.top())], linewidth=0.05)
p += Line(points=[(-2,-2),(-2,-1)], linewidth=0.05, endstyle='>')
p += Line(points=[(2,-2),(2,-1)], linewidth=0.05, endstyle='>')
p += Line(points=[(-2,1),(-2,2),(2,2),(2,1),(2,1),(2,2),(0,2),(G.bottom())], linewidth=0.05, endstyle='>')


p += G
p += P0
p += P1
p += C
print(p)
\end{python}
\begin{console}
#include <iostream>

class G
{};

class P0: public G
{};

class P1: public G
{};

class C: public P0, public P1
{};

int main()
{   
    C c;
    return 0;
}
\end{console}

The important thing to realize is that the object constructions starting
with \verb!C! is done one class at a time going up -- there are
\EMPHASIZE{two paths} in this case. The two paths don't really
``synchronize'' with each other: the creation of objects along the two
inheritance paths is independent. This means that there are \EMPHASIZE{two}
objects created with class G for \verb!c!!!!
\begin{python}
from latextool_basic import *
p = Plot()
#p += Grid(x0=-3,y0=-10,x1=15,y1=5)

G = Rect(-1, 3, 1, 5, linewidth=0.05, font=r'\huge', label=r'\texttt{G}')
P0 = Rect(-3, -1, -0.5, 1, linewidth=0.05, font=r'\huge', label=r'\texttt{P0}')
P1 = Rect(0.5, -1, 3, 1, linewidth=0.05, font=r'\huge', label=r'\texttt{P1}')
C = Rect(-1, -5, 1, -3, linewidth=0.05, font=r'\huge', label=r'\texttt{C}')

p += Line(points=[(-2,-1),(-2,-2),(2,-2),(2,-1),(2,-2),(0,-2),(C.top())], linewidth=0.05)
p += Line(points=[(-2,-2),(-2,-1)], linewidth=0.05, endstyle='>')
p += Line(points=[(2,-2),(2,-1)], linewidth=0.05, endstyle='>')
p += Line(points=[(-2,1),(-2,2),(2,2),(2,1),(2,1),(2,2),(0,2),(G.bottom())], linewidth=0.05, endstyle='>')

E0 = ellipse(5,0,8,2, linewidth=0.1)
E1 = ellipse(9,0,12,2, linewidth=0.1)
E2 = ellipse(5,-3,8,-1, linewidth=0.1)
E3 = ellipse(9,-3,12,-1, linewidth=0.1)
E4 = ellipse(7,-6,10,-4, linewidth=0.1)
p += Line(points=[(6.5,1), (6.5,-2), (8.5,-5), (10.5,-2), (10.5,1)], linewidth=0.1)

p += G
p += Line(points=[G.right(),(6.5,1)], linewidth=0.05)
p += Line(points=[G.right(),(10.5,1)], linewidth=0.05)

p += P0
p += Line(points=[P0.bottomright(),(6.5,-2)], linewidth=0.05)

p += P1
p += Line(points=[P1.right(),(10.5,-2)], linewidth=0.05)

p += C
p += Line(points=[C.right(),(8.5,-5)], linewidth=0.05)


p += E0
p += E1
p += E2
p += E3
p += E4

print(p)
\end{python}

To verify, first do this:
\begin{console}
#include <iostream>

class G
{
public:
        int x_;
};

class P0: public G
{};

class P1: public G
{};

class C: public P0, public P1
{};

int main()
{   
    C c;
    std::cout << &(c.P0::x_) << '\n';
    std::cout << &(c.P1::x_) << '\n';
    return 0;
}
\end{console}
The object \verb!c! has two \verb!x_!'s. We print the address of the
two \verb!x_!'s, which of course are different:

\begin{python}
from latextool_basic import *
p = Plot()
#p += Grid(x0=-3,y0=-10,x1=15,y1=5)

G = Rect(-1, 3, 1, 5, linewidth=0.05, font=r'\huge', label=r'\texttt{G}')
P0 = Rect(-3, -1, -0.5, 1, linewidth=0.05, font=r'\huge', label=r'\texttt{P0}')
P1 = Rect(0.5, -1, 3, 1, linewidth=0.05, font=r'\huge', label=r'\texttt{P1}')
C = Rect(-1, -5, 1, -3, linewidth=0.05, font=r'\huge', label=r'\texttt{C}')

p += Line(points=[(-2,-1),(-2,-2),(2,-2),(2,-1),(2,-2),(0,-2),(C.top())], linewidth=0.05)
p += Line(points=[(-2,-2),(-2,-1)], linewidth=0.05, endstyle='>')
p += Line(points=[(2,-2),(2,-1)], linewidth=0.05, endstyle='>')
p += Line(points=[(-2,1),(-2,2),(2,2),(2,1),(2,1),(2,2),(0,2),(G.bottom())], linewidth=0.05, endstyle='>')

E0 = ellipse(5,0,8,2, linewidth=0.1, label=r'x\_')
E1 = ellipse(9,0,12,2, linewidth=0.1, label=r'x\_')
E2 = ellipse(5,-3,8,-1, linewidth=0.1)
E3 = ellipse(9,-3,12,-1, linewidth=0.1)
E4 = ellipse(7,-6,10,-4, linewidth=0.1)
p += Line(points=[(6.5,1), (6.5,-2), (8.5,-5), (10.5,-2), (10.5,1)], linewidth=0.1)

p += G
p += Line(points=[G.right(),(6.5,1)], linewidth=0.05)
p += Line(points=[G.right(),(10.5,1)], linewidth=0.05)

p += P0
p += Line(points=[P0.bottomright(),(6.5,-2)], linewidth=0.05)

p += P1
p += Line(points=[P1.right(),(10.5,-2)], linewidth=0.05)

p += C
p += Line(points=[C.right(),(8.5,-5)], linewidth=0.05)


p += E0
p += E1
p += E2
p += E3
p += E4

print(p)
\end{python}

\begin{console}
#include <iostream>

class G
{
public:
        int x_;
};

class P0: public G
{};

class P1: public G
{};

class C: public P0, public P1
{};

int main()
{   
    C c;
    std::cout << &(c.P0::x_) << '\n';
    std::cout << &(c.P1::x_) << '\n';
    return 0;
}
\end{console}

Suppose you have a game of shapes. The Shape class describes (with its
subclasses) how to draw the shape. You have another class PointMass that
computes how the shapes moves as a physical point mass. You can have:
\begin{console}
class PhysicalShape: public Shape, PointMass
{...};
\end{console}

This separates out the (motion) physics of object as a point mass from
the drawing of the object.

\begin{itemize}
\item
  If there's a draw method: it's
  probably from the \verb!Shape! class.
\item
  If there's a move method: it's
  probably from the \verb!PointMass! class
\end{itemize}

